% Generated draft: requires mathematical and source review.
\documentclass[11pt]{article}
\usepackage[margin=25mm]{geometry}
\usepackage{fontspec}
\setmainfont{Latin Modern Roman}
\usepackage{amsmath,amssymb,mathtools,unicode-math}
\setmathfont{Latin Modern Math}
\usepackage{enumitem,needspace,xurl,hyperref}
\setlength{\parindent}{0pt}
\setlength{\parskip}{5pt}
\setlength{\emergencystretch}{3em}
\newcommand{\N}{\mathbb{N}}
\newcommand{\Z}{\mathbb{Z}}
\newcommand{\Q}{\mathbb{Q}}
\newcommand{\R}{\mathbb{R}}
\newcommand{\C}{\mathbb{C}}
\newcommand{\F}{\mathbb{F}}
\newcommand{\A}{\mathbb{A}}
\newcommand{\PP}{\mathbb{P}}
\newcommand{\E}{\mathbb{E}}
\newcommand{\eps}{\varepsilon}
\newcommand{\dd}{\,\mathrm{d}}
\begin{document}
\section*{UnsolvedMath ID: 30002454}
\textbf{Source problem code:} catalogue entry
\subsection{Definitions and mathematical statement}
For a field \(K\) and integers \(0\leq k\leq n\), let \(\operatorname{Gr}_K(k,n)\) denote the Grassmannian of \(k\)-dimensional subspaces of \(K^n\). Its Pl\"ucker homogeneous coordinate ring and affine cone are denoted by
\[
R_{k,n,K}=K[p_I: I\subseteq\{1,\ldots,n\},\ |I|=k]/I_{\mathrm{Pl}},
\qquad
X_{k,n,K}=\operatorname{Spec}R_{k,n,K},
\]
where \(I_{\mathrm{Pl}}\) is the ideal of Pl\"ucker relations. These definitions supply notation only: the source record does not specify the intended ground field or the range of \(k,n\) for the question.

\textbf{Source question.} For the affine cone over \(\operatorname{Gr}(k,n)\), viewed through decorated \(\operatorname{GL}(k)\)-connections, do spectral coordinates recover the cluster structure on its homogeneous coordinate ring? In addition, does every cluster in that structure correspond to a spectral network?

More explicitly, if \(\mathscr A\) denotes the intended Grassmannian cluster structure, the second requirement has the schematic quantifiers
\[
\forall\,\mathcal C\in\operatorname{Clust}(\mathscr A),
\quad
\exists\,\mathcal W\text{ a spectral network realizing }\mathcal C.
\]
Here \(\operatorname{Clust}(\mathscr A)\) denotes the set of clusters of \(\mathscr A\). The word ``realizing'' records the requested correspondence; it is not a defined mathematical relation in the supplied record.

\textbf{Unresolved qualifications.} The record does not specify the underlying surface, connection moduli space, decorations, spectral data, admissible networks, or the construction of their spectral coordinates. It also does not specify the identification with the Grassmannian cone, the domains of the coordinate charts, the frozen variables and any localizations, or the equivalence of coordinates intended by ``recover.'' Consequently, the supplied evidence determines this conjectural target but does not determine a fully specified universal proposition. These choices must be recovered from the cited sources before asserting a precise coordinate identification or a correspondence for every cluster.
\subsection{Short English statement}
The question asks whether spectral coordinates associated with decorated \(\operatorname{GL}(k)\)-connections reproduce the Grassmannian cluster structure, and whether every cluster can be realized by a spectral network. The supplied record leaves essential geometric and coordinate conventions unspecified. Its open status is supplied dated curation, not an independently verified conclusion.
\subsection{Sources}
\begin{itemize}
\item Local UnsolvedMath registry: \texttt{data/unsolvedmath/problems.json}, permanent ID \(30002454\), exact source problem code \texttt{catalogue entry}, title ``Spectral Coordinates for Grassmannian Cluster Structures,'' in the Oberwolfach Reports Open Problems collection. This draft uses the supplied complete registry record.
\item Supplied report citation: \emph{Cluster Algebras and Related Topics} (2014), Oberwolfach Reports, pp.~3379--3432, DOI \texttt{10.4171/owr/2013/58}. Original OWR record: \url{https://doi.org/10.4171/owr/2013/58}. The record assigns proposed year 2013 but supplies no individual proposer.
\item Supplied literature source labelled ``Gaiotto Moore Neitzke'': \url{https://arxiv.org/abs/1204.4824}. No further bibliographic details are supplied in the record.
\item Literature assessment checked at \texttt{2026-08-21}: the supplied curation labels the problem open and states that the spectral-coordinate and cluster-coordinate identification remains conjectural, with later examples and partial Grassmannian results rather than a general proof. This is supplied dated curation, not independently verified literature assessment; its time-sensitive status and cited claims require independent checking.
\item The registry labels the statement \texttt{corrected\_verified} and reports comparison with the official Oberwolfach report. That verification claim is part of the supplied provenance and was not independently checked for this draft.
\end{itemize}
\end{document}
